% Options for packages loaded elsewhere
\PassOptionsToPackage{unicode}{hyperref}
\PassOptionsToPackage{hyphens}{url}
\documentclass[
]{article}
\usepackage{xcolor}
\usepackage{amsmath,amssymb}
\setcounter{secnumdepth}{-\maxdimen} % remove section numbering
\usepackage{iftex}
\ifPDFTeX
  \usepackage[T1]{fontenc}
  \usepackage[utf8]{inputenc}
  \usepackage{textcomp} % provide euro and other symbols
\else % if luatex or xetex
  \usepackage{unicode-math} % this also loads fontspec
  \defaultfontfeatures{Scale=MatchLowercase}
  \defaultfontfeatures[\rmfamily]{Ligatures=TeX,Scale=1}
\fi
\usepackage{lmodern}
\ifPDFTeX\else
  % xetex/luatex font selection
\fi
% Use upquote if available, for straight quotes in verbatim environments
\IfFileExists{upquote.sty}{\usepackage{upquote}}{}
\IfFileExists{microtype.sty}{% use microtype if available
  \usepackage[]{microtype}
  \UseMicrotypeSet[protrusion]{basicmath} % disable protrusion for tt fonts
}{}
\makeatletter
\@ifundefined{KOMAClassName}{% if non-KOMA class
  \IfFileExists{parskip.sty}{%
    \usepackage{parskip}
  }{% else
    \setlength{\parindent}{0pt}
    \setlength{\parskip}{6pt plus 2pt minus 1pt}}
}{% if KOMA class
  \KOMAoptions{parskip=half}}
\makeatother
\usepackage{color}
\usepackage{fancyvrb}
\newcommand{\VerbBar}{|}
\newcommand{\VERB}{\Verb[commandchars=\\\{\}]}
\DefineVerbatimEnvironment{Highlighting}{Verbatim}{commandchars=\\\{\}}
% Add ',fontsize=\small' for more characters per line
\newenvironment{Shaded}{}{}
\newcommand{\AlertTok}[1]{\textcolor[rgb]{1.00,0.00,0.00}{\textbf{#1}}}
\newcommand{\AnnotationTok}[1]{\textcolor[rgb]{0.38,0.63,0.69}{\textbf{\textit{#1}}}}
\newcommand{\AttributeTok}[1]{\textcolor[rgb]{0.49,0.56,0.16}{#1}}
\newcommand{\BaseNTok}[1]{\textcolor[rgb]{0.25,0.63,0.44}{#1}}
\newcommand{\BuiltInTok}[1]{\textcolor[rgb]{0.00,0.50,0.00}{#1}}
\newcommand{\CharTok}[1]{\textcolor[rgb]{0.25,0.44,0.63}{#1}}
\newcommand{\CommentTok}[1]{\textcolor[rgb]{0.38,0.63,0.69}{\textit{#1}}}
\newcommand{\CommentVarTok}[1]{\textcolor[rgb]{0.38,0.63,0.69}{\textbf{\textit{#1}}}}
\newcommand{\ConstantTok}[1]{\textcolor[rgb]{0.53,0.00,0.00}{#1}}
\newcommand{\ControlFlowTok}[1]{\textcolor[rgb]{0.00,0.44,0.13}{\textbf{#1}}}
\newcommand{\DataTypeTok}[1]{\textcolor[rgb]{0.56,0.13,0.00}{#1}}
\newcommand{\DecValTok}[1]{\textcolor[rgb]{0.25,0.63,0.44}{#1}}
\newcommand{\DocumentationTok}[1]{\textcolor[rgb]{0.73,0.13,0.13}{\textit{#1}}}
\newcommand{\ErrorTok}[1]{\textcolor[rgb]{1.00,0.00,0.00}{\textbf{#1}}}
\newcommand{\ExtensionTok}[1]{#1}
\newcommand{\FloatTok}[1]{\textcolor[rgb]{0.25,0.63,0.44}{#1}}
\newcommand{\FunctionTok}[1]{\textcolor[rgb]{0.02,0.16,0.49}{#1}}
\newcommand{\ImportTok}[1]{\textcolor[rgb]{0.00,0.50,0.00}{\textbf{#1}}}
\newcommand{\InformationTok}[1]{\textcolor[rgb]{0.38,0.63,0.69}{\textbf{\textit{#1}}}}
\newcommand{\KeywordTok}[1]{\textcolor[rgb]{0.00,0.44,0.13}{\textbf{#1}}}
\newcommand{\NormalTok}[1]{#1}
\newcommand{\OperatorTok}[1]{\textcolor[rgb]{0.40,0.40,0.40}{#1}}
\newcommand{\OtherTok}[1]{\textcolor[rgb]{0.00,0.44,0.13}{#1}}
\newcommand{\PreprocessorTok}[1]{\textcolor[rgb]{0.74,0.48,0.00}{#1}}
\newcommand{\RegionMarkerTok}[1]{#1}
\newcommand{\SpecialCharTok}[1]{\textcolor[rgb]{0.25,0.44,0.63}{#1}}
\newcommand{\SpecialStringTok}[1]{\textcolor[rgb]{0.73,0.40,0.53}{#1}}
\newcommand{\StringTok}[1]{\textcolor[rgb]{0.25,0.44,0.63}{#1}}
\newcommand{\VariableTok}[1]{\textcolor[rgb]{0.10,0.09,0.49}{#1}}
\newcommand{\VerbatimStringTok}[1]{\textcolor[rgb]{0.25,0.44,0.63}{#1}}
\newcommand{\WarningTok}[1]{\textcolor[rgb]{0.38,0.63,0.69}{\textbf{\textit{#1}}}}
\usepackage{longtable,booktabs,array}
\newcounter{none} % for unnumbered tables
\usepackage{calc} % for calculating minipage widths
% Correct order of tables after \paragraph or \subparagraph
\usepackage{etoolbox}
\makeatletter
\patchcmd\longtable{\par}{\if@noskipsec\mbox{}\fi\par}{}{}
\makeatother
% Allow footnotes in longtable head/foot
\IfFileExists{footnotehyper.sty}{\usepackage{footnotehyper}}{\usepackage{footnote}}
\makesavenoteenv{longtable}
\setlength{\emergencystretch}{3em} % prevent overfull lines
\providecommand{\tightlist}{%
  \setlength{\itemsep}{0pt}\setlength{\parskip}{0pt}}
\usepackage{bookmark}
\IfFileExists{xurl.sty}{\usepackage{xurl}}{} % add URL line breaks if available
\urlstyle{same}
\hypersetup{
  hidelinks,
  pdfcreator={LaTeX via pandoc}}

\author{}
\date{}

\begin{document}

\section{Quasar 3.0 --- F-Chain / FHE Red
Review}\label{quasar-30--f-chain--fhe-red-review}

\textbf{Scope}: LP-013 v2 (F-Chain), LP-066 (TFHE), LP-067 (Confidential
ERC-20), LP-068 (Private Teleport), LP-019 (Threshold MPC for keygen),
LP-134 (chain topology), and the corresponding Go reference
(\texttt{lux/precompile/fhe}, \texttt{lux/fhe}).

\textbf{Threat model (as stated)}: adversary controls one F-Chain
validator with full GPU memory visibility, can submit chosen
ciphertexts, can replay/splice messages, cannot break TFHE/Module-LWE.
Goal: leak plaintext or forge state transitions.

\textbf{Reviewer}: Red (adversarial security researcher). \textbf{Date}:
2026-04-26. \textbf{Target launch}: 2025-12-25 (Quasar 3.0 production).
\textbf{Verdict}: \textbf{DO-NOT-SHIP}. The deterministic-keygen, the
unauthenticated \texttt{decrypt} precompile, the lack of
FChainTFHE/FChainBootstrap real verifiers, and the absence of any actual
GPU TFHE kernel mean F-Chain currently provides \textbf{no
confidentiality} and \textbf{no soundness}. Several findings are
showstoppers individually; together they trivially break every
confidential ERC-20 balance and every private teleport.

Numbering: \texttt{Q3.0-FHE-NNN}. Severity: Critical / High / Medium /
Low / Info.

\begin{center}\rule{0.5\linewidth}{0.5pt}\end{center}

\subsection{Executive summary}\label{executive-summary}

{\def\LTcaptype{none} % do not increment counter
\begin{longtable}[]{@{}lll@{}}
\toprule\noalign{}
\# & Title & Severity \\
\midrule\noalign{}
\endhead
\bottomrule\noalign{}
\endlastfoot
Q3.0-FHE-001 & Deterministic keygen → secret key derivable from a public
string & \textbf{Critical} \\
Q3.0-FHE-002 & \texttt{handleDecrypt} has no caller authorization → any
handle decryptable by any contract & \textbf{Critical} \\
Q3.0-FHE-003 & F-Chain cert lanes (\texttt{FChainTFHE},
\texttt{FChainBootstrap}) have no real verifier --- they will inherit
HMAC-keccak placeholder semantics & \textbf{Critical} \\
Q3.0-FHE-004 & None of the GPU TFHE kernels in LP-013 (§Kernels v0.54)
exist on disk & \textbf{Critical} \\
Q3.0-FHE-005 & \texttt{tfheVerify} accepts any byte string that
deserialises --- no input proof, no range check, no type binding &
\textbf{Critical} \\
Q3.0-FHE-006 & \texttt{performFHESealOutput} is a stub: concatenates
pk\textbar ct, returns plaintext-equivalent & \textbf{Critical} \\
Q3.0-FHE-007 & Block-STM treats ciphertexts as opaque bytes → ciphertext
malleability inside an MVCC arena & \textbf{High} \\
Q3.0-FHE-008 & Key-rotation race: two-phase rotation across
\texttt{fchain\_fhe\_root} epochs has no fence; bisecting valid/invalid
ciphertexts leaks the rotation boundary & \textbf{High} \\
Q3.0-FHE-009 & M-Chain → F-Chain bootstrap-key handoff has no published
correctness proof; \texttt{mchain\_ceremony\_root} only commits the
share, not the equivalence to the new \texttt{fchain\_fhe\_root} &
\textbf{High} \\
Q3.0-FHE-010 & F-Chain validator subset → 1/3 corruption halts F-Chain
at lower stake than 1/3 of total Lux & \textbf{High} \\
Q3.0-FHE-011 & \texttt{fhe.decrypt(ciphertext,\ callbackAddress)}
(LP-066) has no nonce: same handle returns same plaintext on repeated
calls --- replay/oracle attacks & \textbf{High} \\
Q3.0-FHE-012 & LP-067 §\texttt{transfer}: no encrypted-amount range
proof tied to caller; underflow-revert leaks the comparison result of
\texttt{amount\ \textgreater{}\ balance} & \textbf{High} \\
Q3.0-FHE-013 & LP-067 transfer pattern (sender, receiver, gas) leaks
even when amount is encrypted; no decoy injection & \textbf{High} \\
Q3.0-FHE-014 & Apple M1/M2 share L2 between SM groups: TFHE bootstrap
timing observable from a co-resident GPU workload & \textbf{Medium} \\
Q3.0-FHE-015 & Three \texttt{external\_product} variants (RGSW × RLWE)
listed without byte-identical correctness gate --- non-deterministic
state root if validators pick different variants & \textbf{Medium} \\
Q3.0-FHE-016 & TLWE noise budget is not tracked anywhere; an adversary
can chain operations to exhaust noise → decryption-error oracle &
\textbf{Medium} \\
Q3.0-FHE-017 & Oblivious branching (\texttt{JUMPI} on encrypted
condition) on \texttt{evm256} is not specified to evaluate both branches
in constant SIMT time --- likely leaks branch-taken via warp divergence
and L1 access pattern & \textbf{Medium} \\
Q3.0-FHE-018 & LP-068 cross-chain replay protection uses
\texttt{nullifier\ =\ Poseidon2(commitment\ ∥\ spending\_key\ ∥\ sourceChainID)}
--- destination chain ID is \textbf{not} in the nullifier; bridge-loop
replay is feasible & \textbf{Medium} \\
Q3.0-FHE-019 & LP-067 storage layout:
\texttt{\_balance{[}address{]}\ =\ ciphertext} per-account encrypted
slot leaks first-time-write vs. update via SSTORE refund pattern (gas
timing oracle) & \textbf{Medium} \\
Q3.0-FHE-020 & Ciphertext handle is \texttt{keccak256(ct\_bytes)} ---
bit-identical ciphertexts collide (two encryptions under the same key
reuse handles in a deterministic-encryption scenario, common with
deterministic seed) & \textbf{Low} \\
Q3.0-FHE-021 & \texttt{tfheRandom} is keyed by user-supplied
\texttt{seed\ uint64} → caller can predict every "random" ciphertext &
\textbf{Low} \\
Q3.0-FHE-022 & \texttt{tfheTrivialEncrypt} uses the encryptor instead of
trivial encryption --- wastes noise budget and confuses the threat model
(LP-013 says trivial constants are noiseless) & \textbf{Info} \\
Q3.0-FHE-023 & LP-067 says transfer reverts when \texttt{sufficient} is
false via underflow --- relies on observable revert, leaking
\texttt{amount\ \textgreater{}\ balance} & (folded into 012) \\
Q3.0-FHE-024 & No PASS evidence anywhere of the 21 GPU kernels claimed
in LP-013 v1; CUDA mirror in v0.55 is also empty & \textbf{Info} \\
Q3.0-FHE-025 & \texttt{evm256} interpreter does not exist in the tree;
LP-013 §"Encrypted EVM (\texttt{evm256.metal})" is design-only &
\textbf{Info} \\
\end{longtable}
}

Severity totals: \textbf{6 Critical, 7 High, 6 Medium, 2 Low, 4 Info}
(25 unique findings).

\begin{center}\rule{0.5\linewidth}{0.5pt}\end{center}

\subsection{Findings}\label{findings}

\subsubsection{Q3.0-FHE-001 --- Deterministic keygen leaks the secret
key to anyone with source
access}\label{q30-fhe-001--deterministic-keygen-leaks-the-secret-key-to-anyone-with-source-access}

\textbf{Severity}: Critical.

\textbf{Files}:

\begin{itemize}
\tightlist
\item
  \texttt{lux/precompile/fhe/fhe\_ops.go:29-60}
\item
  \texttt{lux/fhe/CLAUDE.md} (advertises this fix)
\end{itemize}

\textbf{Scenario}:

\begin{Shaded}
\begin{Highlighting}[]
\KeywordTok{var}\NormalTok{ fheKeygenSeed }\OperatorTok{=}\NormalTok{ sha256}\OperatorTok{.}\NormalTok{Sum256}\OperatorTok{([]}\DataTypeTok{byte}\OperatorTok{(}\StringTok{"LUX\_FHE\_KEYGEN\_v1"}\OperatorTok{))}
\NormalTok{kg}\OperatorTok{,}\NormalTok{ \_ }\OperatorTok{:=}\NormalTok{ fhe}\OperatorTok{.}\NormalTok{NewKeyGeneratorFromSeed}\OperatorTok{(}\NormalTok{params}\OperatorTok{,}\NormalTok{ fheKeygenSeed}\OperatorTok{[:])}
\NormalTok{secretKey}\OperatorTok{,}\NormalTok{ publicKey }\OperatorTok{=}\NormalTok{ kg}\OperatorTok{.}\NormalTok{GenKeyPair}\OperatorTok{()}
\NormalTok{bsk }\OperatorTok{:=}\NormalTok{ kg}\OperatorTok{.}\NormalTok{GenBootstrapKey}\OperatorTok{(}\NormalTok{secretKey}\OperatorTok{)}
\end{Highlighting}
\end{Shaded}

The seed is a public constant. Every node that runs the precompile
derives the same \texttt{(sk,\ pk,\ bsk)}. \textbf{The "secret" key is
universally derivable}, including by the public, not just by validators.
This violates the LP-066 §Security claim "global secret key never exists
in a single location" and the LP-013 §Security claim "only TEE-protected
key holders can decrypt".

\textbf{Attack}: an external observer reads any LP-067 confidential
balance by running \texttt{decryptor.DecryptUint64(ct)} against the
public-constant secret key. There is no attack --- there is no defence.

\textbf{Genesis of the bug}: the prior red swarm (2026-04-13, finding
C-03) flagged a different bug --- \texttt{crypto/rand} keygen made each
validator\textquotesingle s keys diverge. The fix turned the security
parameter into a public constant instead of using a real ceremony
output.

\textbf{Detectability}: zero --- there is no anomaly to monitor.
Compromise is silent and total.

\textbf{Crypto note}: this is not "FHE under a shared key with TEE
custody"; it is "FHE under a key written in the source tree". The
IND-CPA reduction fails immediately:
\(\Pr[\mathcal{A}\text{ wins IND-CPA}] = 1\) for any \(\mathcal{A}\)
that can read the constant.

\textbf{Fix hint}: keys must come from a real ceremony. Either (a)
M-Chain DKG output written into \texttt{fchain\_fhe\_root}, then
validators load shares into TEE on boot and the public key is derived
during the ceremony; or (b) HSM-backed thresholded keygen with
KMS-stored shares (LP-019 §LSS resharing). The seed approach must be
deleted, not parameterised.

\begin{center}\rule{0.5\linewidth}{0.5pt}\end{center}

\subsubsection{\texorpdfstring{Q3.0-FHE-002 --- \texttt{handleDecrypt}
is unauthenticated; any contract can decrypt any
handle}{Q3.0-FHE-002 --- handleDecrypt is unauthenticated; any contract can decrypt any handle}}\label{q30-fhe-002--handledecrypt-is-unauthenticated-any-contract-can-decrypt-any-handle}

\textbf{Severity}: Critical.

\textbf{Files}:

\begin{itemize}
\tightlist
\item
  \texttt{lux/precompile/fhe/contract.go:873-886,\ 1294-1302}
\end{itemize}

\textbf{Scenario}:

\begin{Shaded}
\begin{Highlighting}[]
\KeywordTok{func} \OperatorTok{(}\NormalTok{c }\OperatorTok{*}\NormalTok{FHEContract}\OperatorTok{)}\NormalTok{ handleDecrypt}\OperatorTok{(}\NormalTok{state contract}\OperatorTok{.}\NormalTok{AccessibleState}\OperatorTok{,}\NormalTok{ caller common}\OperatorTok{.}\NormalTok{Address}\OperatorTok{,}\NormalTok{ data }\OperatorTok{[]}\DataTypeTok{byte}\OperatorTok{,}\NormalTok{ gas }\DataTypeTok{uint64}\OperatorTok{,}\NormalTok{ readOnly }\DataTypeTok{bool}\OperatorTok{)} \OperatorTok{(...,} \DataTypeTok{error}\OperatorTok{)} \OperatorTok{\{}
\NormalTok{    handle }\OperatorTok{:=}\NormalTok{ common}\OperatorTok{.}\NormalTok{BytesToHash}\OperatorTok{(}\NormalTok{data}\OperatorTok{[:}\DecValTok{32}\OperatorTok{])}
\NormalTok{    result }\OperatorTok{:=}\NormalTok{ performFHEDecrypt}\OperatorTok{(}\NormalTok{state}\OperatorTok{.}\NormalTok{GetStateDB}\OperatorTok{(),}\NormalTok{ handle}\OperatorTok{,}\NormalTok{ caller}\OperatorTok{)}
    \ControlFlowTok{return}\NormalTok{ result}\OperatorTok{.}\NormalTok{Bytes}\OperatorTok{(),}\NormalTok{ gas }\OperatorTok{{-}}\NormalTok{ GasDecryptRequest}\OperatorTok{,} \OtherTok{nil}
\OperatorTok{\}}

\KeywordTok{func}\NormalTok{ performFHEDecrypt}\OperatorTok{(}\NormalTok{stateDB contract}\OperatorTok{.}\NormalTok{StateDB}\OperatorTok{,}\NormalTok{ handle common}\OperatorTok{.}\NormalTok{Hash}\OperatorTok{,}\NormalTok{ caller common}\OperatorTok{.}\NormalTok{Address}\OperatorTok{)} \OperatorTok{*}\NormalTok{big}\OperatorTok{.}\NormalTok{Int }\OperatorTok{\{}
\NormalTok{    ct}\OperatorTok{,}\NormalTok{ ctType}\OperatorTok{,}\NormalTok{ ok }\OperatorTok{:=}\NormalTok{ getCiphertext}\OperatorTok{(}\NormalTok{stateDB}\OperatorTok{,}\NormalTok{ handle}\OperatorTok{)}
    \ControlFlowTok{if} \OperatorTok{!}\NormalTok{ok }\OperatorTok{\{} \ControlFlowTok{return}\NormalTok{ big}\OperatorTok{.}\NormalTok{NewInt}\OperatorTok{(}\DecValTok{0}\OperatorTok{)} \OperatorTok{\}}
    \ControlFlowTok{return}\NormalTok{ tfheDecrypt}\OperatorTok{(}\NormalTok{ct}\OperatorTok{,}\NormalTok{ ctType}\OperatorTok{)}   \CommentTok{// no ACL, no permit, no callback gating}
\OperatorTok{\}}
\end{Highlighting}
\end{Shaded}

The \texttt{caller} argument is accepted but never consulted. Any
EOA/contract that knows the handle (handles are observable in calldata,
in events, and in state inspection) can call \texttt{decrypt} and get
the plaintext returned in EVM return data. LP-066 §Decryption Protocol
describes a t-of-n threshold decryption with TEE round-trips; the
implementation does single-key decryption synchronously.

\textbf{Attack}:

\begin{enumerate}
\def\labelenumi{\arabic{enumi}.}
\tightlist
\item
  Observe a confidential ERC-20 transfer\textquotesingle s resulting
  \texttt{\_balance{[}victim{]}} handle from the SSTORE trace.
\item
  Craft a contract that calls \texttt{IFHE.decrypt(handle)}.
\item
  Read victim\textquotesingle s balance.
\end{enumerate}

\textbf{Combined with Q3.0-FHE-001}, the precompile
isn\textquotesingle t even needed --- the caller can decrypt off-chain.
But the precompile makes the leak callable from any other contract for
free.

\textbf{Fix hint}: enforce the FheOS ACL pattern
(\texttt{0x0200...0081}) before decryption. The decryption interface
should be split: (a) \texttt{requestDecrypt} emits an event, (b)
off-chain threshold decryptors pick it up, (c) threshold-signed callback
writes the plaintext into the requestor\textquotesingle s state slot via
permit, (d) plaintext never appears in EVM return data.

\begin{center}\rule{0.5\linewidth}{0.5pt}\end{center}

\subsubsection{\texorpdfstring{Q3.0-FHE-003 --- F-Chain cert lanes
(\texttt{FChainTFHE}, \texttt{FChainBootstrap}) inherit HMAC-keccak
placeholder
verification}{Q3.0-FHE-003 --- F-Chain cert lanes (FChainTFHE, FChainBootstrap) inherit HMAC-keccak placeholder verification}}\label{q30-fhe-003--f-chain-cert-lanes-fchaintfhe-fchainbootstrap-inherit-hmac-keccak-placeholder-verification}

\textbf{Severity}: Critical.

\textbf{Files}:

\begin{itemize}
\tightlist
\item
  \texttt{LP-132} §\texttt{drain\_cert\_lane} v0.38 (lines 204-214 of
  \texttt{LP-132-quasar-gpu-execution-adapter.md})
\item
  \texttt{LP-013} §"Cert lanes" (lines 162-168)
\item
  \texttt{LP-134} §\texttt{QuasarCertLane} registry (lines 114-130)
  defines lanes 8 and 9
\end{itemize}

\textbf{Scenario}: LP-132 v0.38 explicitly admits that
\texttt{verify\_bls\_aggregate}, \texttt{verify\_ringtail\_share}, and
\texttt{verify\_mldsa\_groth16} are HMAC-keccak with a master secret.
LP-132 roadmap promises real verifiers at v0.43, v0.44, v0.45.
\textbf{No entry exists for FChainTFHE / FChainBootstrap real
verifiers}. By the LP-013 v0.54 timeline (this LP), F-Chain ships with a
"verifier" that has never been specified. The LP-132 statement
"structured so the swap to real BLS / Ring-LWE / Groth16 is a single
function pointer" does not even consider TFHE attestation, which is
fundamentally a different proof system (proof-of-correct-bootstrap,
e.g., zk-SNARK over a TFHE circuit).

\textbf{Attack}: a single F-Chain validator submits a
\texttt{FChainTFHE} artifact with arbitrary \texttt{circuit\_dag\_root}
and arbitrary plaintext; the HMAC-keccak verifier accepts it because the
validator knows the master secret. The \texttt{fchain\_fhe\_root} no
longer commits to a circuit anyone can verify. Other validators accept
the cert because the round descriptor binds the
artifact\textquotesingle s \texttt{(artifact\_offset,\ artifact\_len)}
indirection --- not the semantic content.

\textbf{Detectability}: the LP-013 §Security claim "Replay protection:
\texttt{fchain\_fhe\_root} binds every cert artifact to the exact
circuit DAG + evaluation-key state" is \textbf{false} under HMAC-keccak:
the verifier never inspects the circuit DAG, only the HMAC tag.

\textbf{Fix hint}: this is the hardest finding. Real F-Chain attestation
requires a SNARK over the TFHE circuit (e.g., verifiable FHE; Lux
already has the \texttt{verifiable-fhe} research scaffold). For Dec 25,
the only viable options are:

\begin{itemize}
\tightlist
\item
  (a) Run F-Chain inside a TEE with attestation through A-Chain
  (\texttt{AChainAttest} lane), constraining trust to TEE compromise; or
\item
  (b) Restrict F-Chain to a permissioned validator set whose stake bond
  is high enough to insure the entire confidential supply (and accept
  that this is a non-cryptographic security argument). Option (a) ducks
  the FHE soundness question by externalising it to TEE. Option (b)
  ducks it by making the attack non-economic. Real verifiable FHE is a
  6-12 month research integration, not a v0.54 patch.
\end{itemize}

\begin{center}\rule{0.5\linewidth}{0.5pt}\end{center}

\subsubsection{Q3.0-FHE-004 --- None of the LP-013 v0.54 GPU TFHE
kernels exist on
disk}\label{q30-fhe-004--none-of-the-lp-013-v054-gpu-tfhe-kernels-exist-on-disk}

\textbf{Severity}: Critical.

\textbf{Files}:

\begin{itemize}
\tightlist
\item
  LP-013 §"Kernels (current as of v0.54)" lists 9 kernels:
  \texttt{tfhe\_bootstrap}, \texttt{blind\_rotate},
  \texttt{blind\_rotate\_fused}, \texttt{external\_product} ×3,
  \texttt{bsk\_prefetch}, \texttt{fhe\_gate}, \texttt{dag\_executor},
  \texttt{evm256}, \texttt{tfhe\_keygen}, \texttt{tfhe\_keyswitch}.
\item
  Searches across \texttt{\textasciitilde{}/work/lux/cevm},
  \texttt{\textasciitilde{}/work/lux/fhe},
  \texttt{\textasciitilde{}/work/lux/precompile/fhe},
  \texttt{\textasciitilde{}/work/lux/luxcpp} (does not exist), and the
  broader tree return \textbf{no matches} for \texttt{tfhe\_bootstrap*},
  \texttt{blind\_rotate*}, \texttt{external\_product*},
  \texttt{drain\_fhe*}, \texttt{fchain\_fhe*}, \texttt{evm256*}. The
  only \texttt{.metal} files on disk are vendored React Native SVG
  filters in \texttt{node\_modules}.
\end{itemize}

\textbf{Implication}: LP-013 v2 is design-only. The Go TFHE reference in
\texttt{lux/fhe} (CPU, \textasciitilde108 ms / bootstrap on M1) is the
entirety of the FHE implementation. The 200 tx/s/GPU on H100 number in
§"Performance targets" is purely aspirational.

\textbf{Attack}: not an attack per se, but a deployment fraud risk:
validators running v0.54 will fall back to CPU TFHE (the LP says GPU is
required for the 500 ms block budget) --- block production halts under
load. The 2026-04-12 fheCRDT audit benchmarked the Go reference at
\textbf{0.24 LWW merges/sec} --- 1,375× slower than what LP-013
promises.

\textbf{Fix hint}: either ship the kernels (months of work --- porting
OpenFHE\textquotesingle s \texttt{bootstrap} to Metal is non-trivial and
the \texttt{cuFHE} reference is GPL) or remove the v0.54 milestone and
admit F-Chain is CPU-bound until real GPU kernels land. \textbf{Do not}
advertise GPU performance numbers in production docs until a kernel
exists.

\begin{center}\rule{0.5\linewidth}{0.5pt}\end{center}

\subsubsection{\texorpdfstring{Q3.0-FHE-005 --- \texttt{tfheVerify} is a
deserialiser, not a
verifier}{Q3.0-FHE-005 --- tfheVerify is a deserialiser, not a verifier}}\label{q30-fhe-005--tfheverify-is-a-deserialiser-not-a-verifier}

\textbf{Severity}: Critical.

\textbf{Files}:

\begin{itemize}
\tightlist
\item
  \texttt{lux/precompile/fhe/fhe\_ops.go:565-568}
\item
  LP-067 §Transfer step 1: "Caller submits an encrypted amount with a ZK
  proof that the ciphertext encrypts a valid uint64."
\end{itemize}

\textbf{Scenario}:

\begin{Shaded}
\begin{Highlighting}[]
\KeywordTok{func}\NormalTok{ tfheVerify}\OperatorTok{(}\NormalTok{ct }\OperatorTok{[]}\DataTypeTok{byte}\OperatorTok{,}\NormalTok{ fheType }\DataTypeTok{uint8}\OperatorTok{)} \DataTypeTok{bool} \OperatorTok{\{}
    \CommentTok{// Basic validation {-} check ciphertext can be deserialized}
    \ControlFlowTok{return}\NormalTok{ deserializeBitCiphertext}\OperatorTok{(}\NormalTok{ct}\OperatorTok{)} \OperatorTok{!=} \OtherTok{nil}
\OperatorTok{\}}
\end{Highlighting}
\end{Shaded}

There is no proof verification. There is no check that the ciphertext
was created with the network public key. There is no check that the
encrypted value lies in the claimed range (e.g.,
\texttt{{[}0,\ 2\^{}64)} for \texttt{euint64}). There is no binding
between ciphertext and caller (otherwise an adversary could re-submit
someone else\textquotesingle s ciphertext as their own).

\textbf{Attack chain (LP-067)}:

\begin{enumerate}
\def\labelenumi{\arabic{enumi}.}
\tightlist
\item
  Adversary submits an encrypted transfer where the "amount" ciphertext
  claims to encrypt 5 (small) but actually encrypts \(2^{64}-1\)
  (giant).
\item
  \texttt{tfheVerify} accepts because deserialisation succeeds.
\item
  Contract computes \texttt{enc\_balance\_sender\ −\ enc\_amount}
  homomorphically.
\item
  Adversary\textquotesingle s balance underflows to
  \textasciitilde{}\(2^{64}\), granting near-infinite confidential token
  balance.
\item
  Contract computes \texttt{enc\_balance\_receiver\ +\ enc\_amount} →
  recipient\textquotesingle s balance also wraps; adversary collaborates
  with a confederate recipient to launder the wrap-around value.
\end{enumerate}

LP-067 names \texttt{inputProof} in the function signature but the
contract trivially passes whatever is supplied to \texttt{tfheVerify},
which discards it.

\textbf{Fix hint}: implement an actual zero-knowledge range proof
(\texttt{R1CS} over the LWE encryption circuit, or simpler: a Sigma
protocol proving the LWE sample\textquotesingle s plaintext is in
\([0, 2^{n})\)). Bind the proof to caller via Fiat-Shamir over
\texttt{(caller,\ ct,\ range,\ public\_key)}. Reject without proof.

\begin{center}\rule{0.5\linewidth}{0.5pt}\end{center}

\subsubsection{\texorpdfstring{Q3.0-FHE-006 ---
\texttt{performFHESealOutput} is a
stub}{Q3.0-FHE-006 --- performFHESealOutput is a stub}}\label{q30-fhe-006--performfhesealoutput-is-a-stub}

\textbf{Severity}: Critical.

\textbf{Files}: \texttt{lux/precompile/fhe/fhe\_ops.go:596-606}.

\textbf{Scenario}:

\begin{Shaded}
\begin{Highlighting}[]
\KeywordTok{func}\NormalTok{ tfheSealOutput}\OperatorTok{(}\NormalTok{ct}\OperatorTok{,}\NormalTok{ pk }\OperatorTok{[]}\DataTypeTok{byte}\OperatorTok{,}\NormalTok{ fheType }\DataTypeTok{uint8}\OperatorTok{)} \OperatorTok{[]}\DataTypeTok{byte} \OperatorTok{\{}
    \CommentTok{// Seal output for a specific public key}
    \CommentTok{// In production, this would re{-}encrypt under the given public key}
    \CommentTok{// For now, just return the ciphertext with a header}
\NormalTok{    result }\OperatorTok{:=} \BuiltInTok{make}\OperatorTok{([]}\DataTypeTok{byte}\OperatorTok{,} \BuiltInTok{len}\OperatorTok{(}\NormalTok{ct}\OperatorTok{)+}\BuiltInTok{len}\OperatorTok{(}\NormalTok{pk}\OperatorTok{)+}\DecValTok{8}\OperatorTok{)}
\NormalTok{    binary}\OperatorTok{.}\NormalTok{BigEndian}\OperatorTok{.}\NormalTok{PutUint32}\OperatorTok{(}\NormalTok{result}\OperatorTok{[}\DecValTok{0}\OperatorTok{:}\DecValTok{4}\OperatorTok{],} \DataTypeTok{uint32}\OperatorTok{(}\BuiltInTok{len}\OperatorTok{(}\NormalTok{pk}\OperatorTok{)))}
\NormalTok{    binary}\OperatorTok{.}\NormalTok{BigEndian}\OperatorTok{.}\NormalTok{PutUint32}\OperatorTok{(}\NormalTok{result}\OperatorTok{[}\DecValTok{4}\OperatorTok{:}\DecValTok{8}\OperatorTok{],} \DataTypeTok{uint32}\OperatorTok{(}\BuiltInTok{len}\OperatorTok{(}\NormalTok{ct}\OperatorTok{)))}
    \BuiltInTok{copy}\OperatorTok{(}\NormalTok{result}\OperatorTok{[}\DecValTok{8}\OperatorTok{:}\DecValTok{8}\OperatorTok{+}\BuiltInTok{len}\OperatorTok{(}\NormalTok{pk}\OperatorTok{)],}\NormalTok{ pk}\OperatorTok{)}
    \BuiltInTok{copy}\OperatorTok{(}\NormalTok{result}\OperatorTok{[}\DecValTok{8}\OperatorTok{+}\BuiltInTok{len}\OperatorTok{(}\NormalTok{pk}\OperatorTok{):],}\NormalTok{ ct}\OperatorTok{)}
    \ControlFlowTok{return}\NormalTok{ result}
\OperatorTok{\}}
\end{Highlighting}
\end{Shaded}

The comment admits it. The "sealed" output is the original ciphertext
plus the requested public key. Under Q3.0-FHE-001 the network secret key
is public so the ciphertext is plaintext-equivalent; under any future
key-ceremony fix, the "seal" is still pointless because no
key-encapsulation has happened --- the recipient gets a ciphertext under
the \textbf{network} key, not their own.

\textbf{Attack}: trivial --- read the bytes after the header.

\textbf{Fix hint}: implement HPKE (\texttt{0x0703} already exists in the
precompile registry) over the recipient\textquotesingle s public key,
encapsulating a fresh AES key that wraps a re-encryption of the
plaintext. Or use proxy re-encryption. Either way: do not ship the stub.

\begin{center}\rule{0.5\linewidth}{0.5pt}\end{center}

\subsubsection{Q3.0-FHE-007 --- Block-STM reads ciphertexts as opaque
bytes; concurrent fiber overwrites are not detected by the homomorphic
circuit}\label{q30-fhe-007--block-stm-reads-ciphertexts-as-opaque-bytes-concurrent-fiber-overwrites-are-not-detected-by-the-homomorphic-circuit}

\textbf{Severity}: High.

\textbf{Files}:

\begin{itemize}
\tightlist
\item
  LP-013 §"Wave-tick co-residency" (lines 128-143)
\item
  LP-010 (QuasarSTM, MVCC arena)
\end{itemize}

\textbf{Scenario}: LP-013 explicitly says "Block-STM sees ciphertext as
opaque bytes for serializability." The MVCC arena\textquotesingle s
RW-set entry records a slot and a version. If two concurrent fibers both
target the same encrypted slot:

\begin{itemize}
\tightlist
\item
  Fiber A: \texttt{SLOAD\ slot\ s\ →\ ct\_old}; computes
  \texttt{ct\_new\ =\ fhe.add(ct\_old,\ x)}; writes back. Read entry:
  \texttt{(s,\ v\_k)}. Write entry: \texttt{(s,\ v\_\{k+1\})}.
\item
  Fiber B (interleaved, before A\textquotesingle s write commits):
  \texttt{SSTORE\ slot\ s\ ←\ ct\_evil} (where \texttt{ct\_evil} is a
  chosen-ciphertext crafted by an adversary).
\item
  Fiber A\textquotesingle s write is detected as a conflict (Block-STM
  repair) and re-runs with \texttt{ct\_evil} as \texttt{ct\_old}. The
  homomorphic circuit happily computes \texttt{fhe.add(ct\_evil,\ x)} →
  state slot now contains \texttt{ct\_evil\ +\ x}.
\end{itemize}

The homomorphic gate has no way to detect that \texttt{ct\_evil} is
malformed (see Q3.0-FHE-005). Block-STM only enforces
\emph{serialisability} of writes, not \emph{semantic validity} of
ciphertexts.

\textbf{Attack}: classic confidential ERC-20 "front-run with a rogue
ciphertext" --- adversary submits a tx that writes \texttt{ct\_evil} (a
ciphertext of a giant value crafted to cause a specific underflow when
added to the victim\textquotesingle s balance). Block-STM serialises it
before the victim\textquotesingle s transfer. Victim\textquotesingle s
transfer then operates on \texttt{ct\_evil}, propagating the corruption
into the victim\textquotesingle s balance.

\textbf{Fix hint}: every encrypted SSTORE must validate the ciphertext
via the range-proof gate from Q3.0-FHE-005 \emph{before} the MVCC slot
accepts the write. The range proof must be tied to the writing caller,
not the reader.

\begin{center}\rule{0.5\linewidth}{0.5pt}\end{center}

\subsubsection{\texorpdfstring{Q3.0-FHE-008 --- Key-rotation race across
\texttt{fchain\_fhe\_root} epochs has no
fence}{Q3.0-FHE-008 --- Key-rotation race across fchain\_fhe\_root epochs has no fence}}\label{q30-fhe-008--key-rotation-race-across-fchain_fhe_root-epochs-has-no-fence}

\textbf{Severity}: High.

\textbf{Files}:

\begin{itemize}
\tightlist
\item
  LP-013 §"Key management" (lines 171-184)
\end{itemize}

\textbf{Scenario}: LP-013 says "key rotation is an M-Chain ceremony that
emits a new \texttt{fchain\_fhe\_root}; F-Chain finalizes by accepting a
Quasar round whose descriptor cites the new root." The spec does not
define:

\begin{itemize}
\tightlist
\item
  The atomicity boundary (a single Quasar round, or a window?)
\item
  Whether txs in flight at rotation use old key or new key
\item
  What happens to ciphertexts under the old key after rotation
  (re-encryption? grandfathering? deletion?)
\end{itemize}

\textbf{Attack (rotation oracle)}: an adversary submits two
near-identical transfer ciphertexts \(ct_0\) (under old key) and
\(ct_1\) (under new key) back-to-back. Whichever succeeds tells the
adversary the rotation epoch boundary precisely. By repeating across
many slots, the adversary builds a time series of rotation events ---
useful for correlating off-chain TEE maintenance windows, validator
restart patterns, and KMS rotation schedules.

\textbf{Attack (split-brain)}: if some fibers in the same wave-tick use
the old \texttt{fchain\_fhe\_root} and some use the new (because the
descriptor was swapped mid-tick by a malicious miner submitting two
competing descriptors), the resulting state root is non-deterministic
across validators. Either a \texttt{repair} round restores consistency
(liveness loss) or validators sign different roots (safety loss).

\textbf{Fix hint}: the ceremony must produce two outputs: (1) the new
\texttt{fchain\_fhe\_root}, (2) a signed cross-key re-encryption proof
that maps every active ciphertext slot to its new-key equivalent.
Rotation lands atomically as a Quasar round whose descriptor binds both.
Until the re-encryption is committed, the new key cannot decrypt and the
old key must not be deleted. Block this as a \texttt{HotLaneMode} to
serialise.

\begin{center}\rule{0.5\linewidth}{0.5pt}\end{center}

\subsubsection{Q3.0-FHE-009 --- M-Chain → F-Chain handoff lacks
correctness
binding}\label{q30-fhe-009--m-chain--f-chain-handoff-lacks-correctness-binding}

\textbf{Severity}: High.

\textbf{Files}:

\begin{itemize}
\tightlist
\item
  LP-013 §"Key management" (line 174)
\item
  LP-019 §CGGMP21 / FROST (general MPC, no TFHE-specific path)
\item
  LP-066 §"Key Generation" (mentions TEE mesh DKG but no LP for it)
\end{itemize}

\textbf{Scenario}: LP-013 says "TFHE keys (bootstrap key, key-switch
key) are generated via M-Chain MPC ceremony (LP-019 §TFHE keygen,
LP-076)." LP-019 has no §TFHE keygen. LP-066 says "global FHE public key
is generated via a distributed key generation ceremony among TEE mesh
participants (LP-065)" --- a different ceremony, not in scope of
M-Chain.

The chain \texttt{M\ →\ F} of trust is therefore unspecified:

\begin{itemize}
\tightlist
\item
  Who runs the TFHE DKG? CGGMP21 produces ECDSA shares, not TFHE shares.
  TFHE thresholding requires a different DKG (e.g., TFHE Threshold from
  Mouchet, ASIACRYPT\textquotesingle22).
\item
  What does \texttt{mchain\_ceremony\_root} commit to? If it commits to
  a CGGMP21 group public key, that does not constrain a TFHE bootstrap
  key.
\item
  The \texttt{MChainCGGMP21} cert lane only proves a CGGMP21 share is
  valid; it does not prove that a downstream TFHE key was constructed
  honestly from it.
\end{itemize}

\textbf{Attack}: a malicious M-Chain coordinator forges a
\texttt{mchain\_ceremony\_root} that commits to a TFHE bootstrap key
generated by them alone (not via DKG). They send the malicious key to
F-Chain, which has no way to verify that the key is the output of a
t-of-n ceremony --- \texttt{fchain\_fhe\_root} just hashes whatever
F-Chain accepts. Now the malicious party knows the secret key for all
confidential ERC-20 balances.

\textbf{Fix hint}: define an explicit TFHE DKG (LP-019 needs a §TFHE-DKG
section). The DKG must produce a verifiable transcript; F-Chain must
accept the new \texttt{fchain\_fhe\_root} only when the transcript hash
is bound to \texttt{mchain\_ceremony\_root} and the transcript is
publicly auditable. Use an existing scheme (Mouchet et al., or
Boneh-Gentry-Halevi-Wang). Do not ship without this LP.

\begin{center}\rule{0.5\linewidth}{0.5pt}\end{center}

\subsubsection{Q3.0-FHE-010 --- F-Chain validator subset reduces
1/3-stake
threshold}\label{q30-fhe-010--f-chain-validator-subset-reduces-13-stake-threshold}

\textbf{Severity}: High.

\textbf{Files}:

\begin{itemize}
\tightlist
\item
  LP-134 §F-Chain (lines 276-288)
\item
  LP-134 §"VM identifiers" (\texttt{F\ =\ lux:fhe})
\end{itemize}

\textbf{Scenario}: F-Chain has its own validator set carved from
P-Chain. The spec is silent on the size of the F-Chain validator set
relative to the total Lux validator set. Quasar BFT tolerance is 1/3 of
\emph{that subset\textquotesingle s} stake, not of total Lux stake.

\textbf{Attack}:

\begin{itemize}
\tightlist
\item
  If F-Chain runs with, say, 21 validators carved from a 100-validator
  Lux mainnet, an adversary needs ≥ 8 of 21 = \textasciitilde38\% of
  F-Chain stake to halt F-Chain. That may be only \textasciitilde8\% of
  total Lux stake. \textbf{The cost of attacking F-Chain is 4× lower
  than the cost of attacking the rest of the network.}
\item
  For confidentiality (Q3.0-FHE-009), if the F-Chain validator set is
  also the M-Chain TFHE-DKG ceremony set, the same adversary can corrupt
  the DKG with ≥ 1/3 of F-Chain stake (regardless of Quasar 2/3 BFT,
  because DKG only requires t honest parties out of n).
\end{itemize}

\textbf{Fix hint}: either (a) require F-Chain validators to be a strict
super-set of P-Chain top-stake validators (e.g., top 67\% by stake), so
F-Chain compromise implies network compromise; (b) require F-Chain
validators to post additional bond proportional to total confidential
supply on the chain; (c) explicitly publish the F-Chain set and its
stake distribution and let dApps decide whether to trust it.

\begin{center}\rule{0.5\linewidth}{0.5pt}\end{center}

\subsubsection{Q3.0-FHE-011 --- Synchronous decrypt is a deterministic
oracle}\label{q30-fhe-011--synchronous-decrypt-is-a-deterministic-oracle}

\textbf{Severity}: High.

\textbf{Files}:

\begin{itemize}
\tightlist
\item
  LP-066 §"Decryption Protocol": "When t partial decryptions are
  collected, the result is aggregated and delivered to the callback.
  Decryption latency: \textasciitilde2 seconds."
\item
  \texttt{lux/precompile/fhe/contract.go:873-886} --- implementation is
  \textbf{synchronous}, no callback.
\end{itemize}

\textbf{Scenario}: LP-066 specifies async threshold decryption. The
implementation does single-key sync decryption, returning the plaintext
in EVM return data with no nonce, no timestamp, no caller binding.
Repeated calls on the same handle return the same plaintext
(deterministic). This is a textbook decryption oracle.

\textbf{Attack (with Q3.0-FHE-002)}:

\begin{enumerate}
\def\labelenumi{\arabic{enumi}.}
\tightlist
\item
  Adversary calls \texttt{decrypt(handle\_v)} and gets
  victim\textquotesingle s plaintext value.
\item
  If victim later updates the slot with a different ciphertext under the
  same key (e.g., same plaintext but fresh randomness), the adversary
  can correlate the new handle to the old plaintext via known-plaintext
  arithmetic:
  \texttt{decrypt(fhe.sub(new\_handle,\ old\_handle))\ =\ delta}.
\item
  The chain of deltas reveals the entire balance trajectory.
\end{enumerate}

\textbf{Fix hint}: remove the sync \texttt{decrypt} precompile entirely.
Implement \texttt{requestDecrypt(handle)\ →\ (requestId,\ event)};
off-chain threshold decryptors deliver via signed callback; permit-bound
delivery into a caller-only state slot.

\begin{center}\rule{0.5\linewidth}{0.5pt}\end{center}

\subsubsection{Q3.0-FHE-012 --- LP-067 underflow-revert pattern leaks
comparison
result}\label{q30-fhe-012--lp-067-underflow-revert-pattern-leaks-comparison-result}

\textbf{Severity}: High.

\textbf{Files}:

\begin{itemize}
\tightlist
\item
  LP-067 §Transfer step 5: "If \texttt{sufficient} is false, the
  subtraction underflows and the transaction reverts."
\end{itemize}

\textbf{Scenario}: the contract intentionally lets underflow signal
insufficient balance via a revert. Whether the tx reverts is publicly
observable on-chain. \textbf{Therefore the comparison
\texttt{amount\ \textgreater{}\ balance(sender)} is publicly observable}
--- even though both operands are encrypted. This trivially leaks bits
about the sender\textquotesingle s balance.

\textbf{Attack}:

\begin{enumerate}
\def\labelenumi{\arabic{enumi}.}
\tightlist
\item
  Adversary controls a faucet contract that lets them submit a transfer
  with a chosen plaintext amount \(a\).
\item
  Adversary issues transfers of amount \(a_1, a_2, \ldots\) to a probe
  address, observing which revert and which succeed.
\item
  Binary search recovers victim\textquotesingle s balance to arbitrary
  precision in \(O(\log_2 \text{balance})\) transfers.
\end{enumerate}

\textbf{Fix hint}: the pattern must be \texttt{cmux} over the
comparison. If \texttt{sufficient\ =\ FHE.le(amount,\ balance)}, then:

\begin{itemize}
\tightlist
\item
  \texttt{new\_sender\ =\ FHE.sub(balance,\ FHE.cmux(sufficient,\ amount,\ 0))}
\item
  \texttt{new\_recv\ \ \ =\ FHE.add(recv,\ \ \ \ FHE.cmux(sufficient,\ amount,\ 0))}
\end{itemize}

Tx never reverts due to balance; instead it commits a no-op when
insufficient. Then the \emph{fact} of insufficiency is itself encrypted
(the recipient simply does not receive anything --- observable but
cleartext- indistinguishable from a legitimate zero-amount transfer if
the caller issues such transfers as decoys). Combined with batched
transfers and nonce-gated visibility, the comparison oracle closes.

\begin{center}\rule{0.5\linewidth}{0.5pt}\end{center}

\subsubsection{Q3.0-FHE-013 --- LP-067 leaks transfer graph (sender,
receiver,
gas)}\label{q30-fhe-013--lp-067-leaks-transfer-graph-sender-receiver-gas}

\textbf{Severity}: High.

\textbf{Files}:

\begin{itemize}
\tightlist
\item
  LP-067 §"Security Considerations" item 3 already admits "metadata
  leakage: sender, receiver, and timing are still visible".
\end{itemize}

\textbf{Scenario}: even with perfect amount confidentiality,
transfer-graph analysis reveals account clusters, employer--employee
relationships, DEX arbitrage flows, etc. LP-067 punts to LP-068, but
LP-068 only addresses \emph{cross-chain} private transfers, not
\emph{intra-chain}. There is no privacy mixer for confidential ERC-20 in
scope.

\textbf{Additional gas leak}: TFHE operations cost gas proportional to
the number of bootstraps, which depends on the bit-width and the
operation. Different ciphertext sizes (e.g., euint8 vs. euint64) reveal
the \emph{type} of transfer. Even within euint64, the gas usage for the
first SSTORE into a fresh slot is higher than an update SSTORE --- this
distinguishes "first incoming transfer to a wallet" from "subsequent
transfer".

\textbf{Attack}: passive on-chain analysis. Standard graph-mining tools
(used on Bitcoin / Ethereum) work unchanged because the FHE only
encrypts amounts, not addresses or gas.

\textbf{Fix hint}:

\begin{itemize}
\tightlist
\item
  Combine LP-067 with a privacy mixer (LP-064 ShieldedPool) so that
  amounts AND addresses are confidential. LP-067 alone is not a privacy
  primitive; it is a confidentiality primitive.
\item
  Pad SSTORE gas to a constant cost regardless of slot freshness (use a
  dedicated FHE storage subspace where every slot is pre-initialised
  with an encrypted zero at deploy time).
\end{itemize}

\begin{center}\rule{0.5\linewidth}{0.5pt}\end{center}

\subsubsection{Q3.0-FHE-014 --- Apple M1/M2 GPU L2 cache shared between
SM
groups}\label{q30-fhe-014--apple-m1m2-gpu-l2-cache-shared-between-sm-groups}

\textbf{Severity}: Medium.

\textbf{Files}: LP-013 §"Performance targets" (Apple M1 Max numbers).

\textbf{Scenario}: Apple Silicon GPUs share L2 cache between SM groups
(each SM has its own L1 but contends for L2). TFHE bootstrap operates on
the RGSW × RLWE external product, which hits L2 hard for the BSK. A
co-resident GPU workload (e.g., a malicious WebGPU page running in
Safari, or another Metal compute task on a multi-tenant macOS validator)
can:

\begin{itemize}
\tightlist
\item
  Probe L2 latency (PRIME+PROBE) to learn which BSK rows are hot in a
  given bootstrap.
\item
  Correlate hot rows across many bootstraps to recover the secret-key
  decomposition (BSK = bit-decomposed encryption of secret key under
  RGSW; cache hits leak the secret-key bits).
\end{itemize}

This is not science fiction: similar attacks against AES (Bernstein),
RSA (FLUSH+RELOAD), and AES-NI on Intel are well-established. Apple has
no documented L2 partitioning primitive.

\textbf{Attack complexity}: high (requires co-residency on a validator
GPU) but feasible if a validator runs other GPU workloads (LLM
inference, gaming, browser GPU, HPL) on the same M1 Max alongside
\texttt{drain\_fhe}.

\textbf{Fix hint}: F-Chain validators must run on dedicated GPUs with no
other workload. Document this as a hard requirement. Better: use H100
with MIG partitioning, where L2 is partitioned by hardware. On Apple
Silicon, validators must dedicate the entire GPU to \texttt{drain\_fhe}
and disable Safari/Metal preemption. Until LP-013 specifies this, treat
Apple Silicon F-Chain validators as best-effort (testnet only).

\begin{center}\rule{0.5\linewidth}{0.5pt}\end{center}

\subsubsection{\texorpdfstring{Q3.0-FHE-015 --- Three
\texttt{external\_product} variants without correctness
gate}{Q3.0-FHE-015 --- Three external\_product variants without correctness gate}}\label{q30-fhe-015--three-external_product-variants-without-correctness-gate}

\textbf{Severity}: Medium.

\textbf{Files}: LP-013 §"Kernels" ---
\texttt{external\_product\ (×3\ variants)}.

\textbf{Scenario}: the spec lists three variants of RGSW × RLWE external
product but does not specify a byte-equality requirement on outputs.
Variants typically differ in NTT layout, decomposition base, or BSK
prefetch strategy. If two variants produce \emph{equivalent} but not
\emph{byte-identical} RLWE samples (e.g., differing in noise term while
decrypting to the same plaintext), then validators using different
variants will produce different \texttt{state\_root} values --- fork.

\textbf{Attack}: a malicious miner crafts a tx that exercises a code
path where Variant A and Variant B diverge. Validators with Variant A
produce state root \(r_A\); validators with Variant B produce \(r_B\).
The chain forks. The miner times the attack to a low-quorum window to
maximise chain instability.

\textbf{Fix hint}: pick one variant. Period. The three-variant
flexibility is a performance optimisation that contradicts deterministic
execution. LP-013 must specify \emph{the} variant (with parameters) and
ban the others on the cert path. If a variant is dropped later, gate it
behind a hard fork.

\begin{center}\rule{0.5\linewidth}{0.5pt}\end{center}

\subsubsection{Q3.0-FHE-016 --- TLWE noise budget is not tracked
anywhere}\label{q30-fhe-016--tlwe-noise-budget-is-not-tracked-anywhere}

\textbf{Severity}: Medium.

\textbf{Files}:

\begin{itemize}
\tightlist
\item
  \texttt{lux/fhe/evaluator.go} (no noise tracking)
\item
  LP-013 §Security: "TFHE gates run in constant time per bootstrap"
  (true but unrelated)
\item
  2026-04-12 fhecrdt-audit §"No noise budget monitoring" (already
  flagged)
\end{itemize}

\textbf{Scenario}: TFHE bootstrap resets noise to a known level \emph{if
applied correctly}. But:

\begin{itemize}
\tightlist
\item
  The Go reference uses \texttt{eval.bootstrap} inside every gate, which
  is fine for boolean gates. For arithmetic ops (Add, Mul) over
  \texttt{BitCiphertext} (multi-bit), the bootstrap is applied per-bit.
  If a downstream caller treats a \texttt{BitCiphertext} as a unit and
  chains operations without bootstrap (e.g., scalar multiplication
  followed by addition), noise can accumulate.
\item
  The chain of \texttt{fhe.add(fhe.add(fhe.add(...)))} in confidential
  ERC-20 --- each \texttt{Add} does internal bootstrap, so this
  \emph{should} be safe. But unverified.
\item
  An adversary can construct a tx that intentionally exhausts noise via
  an unbootstrapped chain (if any exists), causing decryption to flip
  bits with attacker-chosen probability --- a key-recovery oracle
  (Loftus- May-Smart-Vercauteren).
\end{itemize}

\textbf{Attack}: locate any unbootstrapped chain in the gate library;
submit a contract that drives it; observe the decryption-error pattern
via the Q3.0-FHE-002 oracle; recover the secret key (if Q3.0-FHE-001
weren\textquotesingle t already trivial).

\textbf{Fix hint}: add a property test that asserts every public TFHE op
(\texttt{Add}, \texttt{Sub}, \texttt{Mul}, \texttt{Div}, \texttt{Rem},
\texttt{Lt}, \texttt{Le}, \texttt{Eq}, \texttt{Cmux}, \texttt{Min},
\texttt{Max}) leaves the output noise at exactly the post-bootstrap
level. Run in CI on every commit. Reject ops that fail the property.

\begin{center}\rule{0.5\linewidth}{0.5pt}\end{center}

\subsubsection{\texorpdfstring{Q3.0-FHE-017 --- Oblivious branching not
specified in
\texttt{evm256}}{Q3.0-FHE-017 --- Oblivious branching not specified in evm256}}\label{q30-fhe-017--oblivious-branching-not-specified-in-evm256}

\textbf{Severity}: Medium.

\textbf{Files}: LP-013 §"Encrypted EVM (\texttt{evm256.metal})" ---
"Branching (JUMPI on encrypted condition) via fully oblivious evaluation
of both branches, then ciphertext-mux."

\textbf{Scenario}: "fully oblivious" requires:

\begin{enumerate}
\def\labelenumi{\arabic{enumi}.}
\tightlist
\item
  Both branches execute regardless of condition.
\item
  Both branches consume the \textbf{same} number of gates / cycles /
  cache lines / memory accesses.
\item
  The mux is applied as a homomorphic gate, not a CPU select.
\end{enumerate}

In SIMT (Apple Metal, NVIDIA CUDA), warp divergence on an encrypted
condition is not directly possible (the condition is a ciphertext bit,
not a hardware predicate), but \textbf{memory access patterns} can still
diverge if the two branches touch different state slots. The Block-STM
read set will differ, leaking which branch was "taken" via the RW set.

\textbf{Attack}: an adversary writes a contract:

\begin{Shaded}
\begin{Highlighting}[]
\NormalTok{ebool cond = FHE.lt(victim\_balance, threshold);}
\NormalTok{if (cond) sstore(slot\_A, ...);}
\NormalTok{else      sstore(slot\_B, ...);}
\end{Highlighting}
\end{Shaded}

After tx execution, the read/write set in the round descriptor reveals
whether \texttt{slot\_A} or \texttt{slot\_B} was written, leaking
\texttt{victim\_balance\ \textless{}\ threshold}.

\textbf{Fix hint}: the contract compiler must lower encrypted JUMPI to:

\begin{Shaded}
\begin{Highlighting}[]
\NormalTok{v\_A = compute\_branch\_A();}
\NormalTok{v\_B = compute\_branch\_B();}
\NormalTok{sstore(slot\_A, FHE.cmux(cond, v\_A, sload(slot\_A)));}
\NormalTok{sstore(slot\_B, FHE.cmux(cond, sload(slot\_B), v\_B));}
\end{Highlighting}
\end{Shaded}

i.e., always read+write both slots, mux the result. The RW set then
includes both slots regardless of condition. This is expensive
(\textasciitilde2× gas) but is the only correct path for oblivious
branching at the EVM level.

\begin{center}\rule{0.5\linewidth}{0.5pt}\end{center}

\subsubsection{Q3.0-FHE-018 --- LP-068 cross-chain replay
protection}\label{q30-fhe-018--lp-068-cross-chain-replay-protection}

\textbf{Severity}: Medium.

\textbf{Files}:

\begin{itemize}
\tightlist
\item
  LP-068 §"Cross-Chain Nullifier Sync" (line 56):
  \texttt{nullifier\ =\ Poseidon2(commitment\ ∥\ spending\_key\ ∥\ sourceChainID)}
\end{itemize}

\textbf{Scenario}: the destination chain ID is \textbf{not} part of the
nullifier. A user deposits on chain \texttt{S}, withdraws on chain
\texttt{D1}. The nullifier is \(N = H(\text{commitment} \| sk \| S)\),
registered on \(D1\). Now the warp message also reaches \(D2\) (or an
attacker re-routes it). \(D2\) has not seen \(N\), so the withdrawal
succeeds again. The attacker drains liquidity on \(D2\).

LP-068 §Security item 4 claims "Cross-chain replay is prevented by
including the source chain ID in the nullifier" --- but this prevents
\emph{source-side} replay, not \emph{destination-side} replay. The
destination chain ID being absent from the nullifier means the nullifier
is the same on every destination chain. A relayer can re-broadcast the
same warp message to multiple destinations.

\textbf{Fix hint}: include destination in the nullifier:
\(N = H(\text{commitment} \| sk \| S \| D \| \text{sequence})\).
Nullifier is then unique per source-destination pair. Maintain a global
nullifier registry (B-Chain) that aggregates nullifiers across all
chains; reject any nullifier already seen on any destination.

\begin{center}\rule{0.5\linewidth}{0.5pt}\end{center}

\subsubsection{Q3.0-FHE-019 --- SSTORE refund pattern leaks first-write
vs.
update}\label{q30-fhe-019--sstore-refund-pattern-leaks-first-write-vs-update}

\textbf{Severity}: Medium.

\textbf{Files}: \texttt{lux/precompile/fhe/contract.go:920-994}
(storeCiphertext loop).

\textbf{Scenario}: the \texttt{storeCiphertext} function writes the
ciphertext as a sequence of 32-byte chunks via
\texttt{stateDB.SetState}. EVM SSTORE has different gas costs for fresh
slots (20 000 gas) vs. update slots (5 000 gas) vs. zero-clear (refund
15 000 gas). The total gas for a confidential ERC-20 transfer therefore
reveals:

\begin{itemize}
\tightlist
\item
  Whether the sender\textquotesingle s balance slot already existed
  (existing user vs. new user).
\item
  Whether the recipient\textquotesingle s balance slot already existed.
\item
  The size class of the ciphertext (uint8 = 8 KB = 256 chunks vs.
  uint256 = 32 KB = 1024 chunks).
\end{itemize}

\textbf{Attack}: passive gas analysis. Surveillance is straightforward.

\textbf{Fix hint}: pre-allocate maximum-size slots at deployment time;
overwrite in place. Pad ciphertexts to a fixed size class (e.g., always
32 KB). Use SSTORE patterns that have constant gas cost (write 32 KB of
zeros first, then overwrite with ciphertext --- accept the extra gas as
the cost of confidentiality).

\begin{center}\rule{0.5\linewidth}{0.5pt}\end{center}

\subsubsection{\texorpdfstring{Q3.0-FHE-020 --- Handle =
\texttt{keccak256(ct)} collides on bit-identical
ciphertexts}{Q3.0-FHE-020 --- Handle = keccak256(ct) collides on bit-identical ciphertexts}}\label{q30-fhe-020--handle--keccak256ct-collides-on-bit-identical-ciphertexts}

\textbf{Severity}: Low.

\textbf{Files}: \texttt{lux/precompile/fhe/contract.go:935-936}.

\textbf{Scenario}: \texttt{handle\ =\ keccak256(ct)}. Two ciphertexts
with identical bytes produce the same handle. Under randomised LWE
encryption, identical ciphertexts are negligible probability; but with
the deterministic seed in Q3.0-FHE-001, a deterministic-encryption mode
(e.g., \texttt{tfheTrivialEncrypt(constant)}) collides for identical
plaintexts. Two separate users encrypting "10" trivially get the same
handle; the second user\textquotesingle s encryption silently aliases to
the first\textquotesingle s slot.

\textbf{Fix hint}: handle should include caller and a sequence:
\(H = \text{keccak256}(caller \| nonce \| ct)\).

\begin{center}\rule{0.5\linewidth}{0.5pt}\end{center}

\subsubsection{\texorpdfstring{Q3.0-FHE-021 --- \texttt{tfheRandom} is
keyed by user-supplied
seed}{Q3.0-FHE-021 --- tfheRandom is keyed by user-supplied seed}}\label{q30-fhe-021--tfherandom-is-keyed-by-user-supplied-seed}

\textbf{Severity}: Low.

\textbf{Files}: \texttt{lux/precompile/fhe/fhe\_ops.go:608-619} (called
via \texttt{handleRand}).

\textbf{Scenario}:

\begin{Shaded}
\begin{Highlighting}[]
\KeywordTok{func}\NormalTok{ tfheRandom}\OperatorTok{(}\NormalTok{fheType }\DataTypeTok{uint8}\OperatorTok{,}\NormalTok{ seed }\DataTypeTok{uint64}\OperatorTok{)} \OperatorTok{[]}\DataTypeTok{byte} \OperatorTok{\{}
\NormalTok{    seedBytes }\OperatorTok{:=} \BuiltInTok{make}\OperatorTok{([]}\DataTypeTok{byte}\OperatorTok{,} \DecValTok{32}\OperatorTok{)}
\NormalTok{    binary}\OperatorTok{.}\NormalTok{BigEndian}\OperatorTok{.}\NormalTok{PutUint64}\OperatorTok{(}\NormalTok{seedBytes}\OperatorTok{[}\DecValTok{24}\OperatorTok{:],}\NormalTok{ seed}\OperatorTok{)}
\NormalTok{    rng }\OperatorTok{:=}\NormalTok{ fhe}\OperatorTok{.}\NormalTok{NewFheRNG}\OperatorTok{(}\NormalTok{params}\OperatorTok{,}\NormalTok{ secretKey}\OperatorTok{,}\NormalTok{ seedBytes}\OperatorTok{)}
    \CommentTok{// ...}
\OperatorTok{\}}
\end{Highlighting}
\end{Shaded}

The caller supplies the seed. There is no entropy beyond the
caller\textquotesingle s input. Two callers supplying the same seed
produce the same encrypted "random" value. A confidential gambling
contract that reads \texttt{fhe.rand} is trivially predictable.

Worse: the function uses \texttt{secretKey} (Q3.0-FHE-001 → public) plus
the caller\textquotesingle s seed → the "random" output is
deterministically computable off-chain by anyone.

\textbf{Fix hint}: derive the seed from
\texttt{(blockhash,\ tx\_hash,\ salt)} and HKDF through the threshold
key. Or use a VRF (LP-019 already has FROST-VRF). Never accept
caller-controlled seeds for on-chain randomness.

\begin{center}\rule{0.5\linewidth}{0.5pt}\end{center}

\subsubsection{\texorpdfstring{Q3.0-FHE-022 ---
\texttt{tfheTrivialEncrypt} is not trivial
encryption}{Q3.0-FHE-022 --- tfheTrivialEncrypt is not trivial encryption}}\label{q30-fhe-022--tfhetrivialencrypt-is-not-trivial-encryption}

\textbf{Severity}: Info.

\textbf{Files}: \texttt{lux/precompile/fhe/fhe\_ops.go:584-594}.

\textbf{Scenario}:

\begin{Shaded}
\begin{Highlighting}[]
\KeywordTok{func}\NormalTok{ tfheTrivialEncrypt}\OperatorTok{(}\NormalTok{plaintext }\OperatorTok{*}\NormalTok{big}\OperatorTok{.}\NormalTok{Int}\OperatorTok{,}\NormalTok{ toType }\DataTypeTok{uint8}\OperatorTok{)} \OperatorTok{[]}\DataTypeTok{byte} \OperatorTok{\{}
    \OperatorTok{...}
\NormalTok{    ct }\OperatorTok{:=}\NormalTok{ encryptor}\OperatorTok{.}\NormalTok{EncryptUint64}\OperatorTok{(}\NormalTok{plaintext}\OperatorTok{.}\NormalTok{Uint64}\OperatorTok{(),}\NormalTok{ targetType}\OperatorTok{)}
    \ControlFlowTok{return}\NormalTok{ serializeBitCiphertext}\OperatorTok{(}\NormalTok{ct}\OperatorTok{)}
\OperatorTok{\}}
\end{Highlighting}
\end{Shaded}

Comment says "Use encryptor for now (trivial encryption would be
noiseless)". Trivial encryption in TFHE is "embed the plaintext directly
with zero noise" --- which would be safe for public constants. The
current code uses real encryption with real noise, wasting the noise
budget for what should be a freebie. Not a security bug, but a
correctness debt.

\textbf{Fix hint}: implement actual trivial encryption (noiseless RLWE
sample with the plaintext in the constant term). Document that trivial
encryption MUST NOT be used for secret data.

\begin{center}\rule{0.5\linewidth}{0.5pt}\end{center}

\subsubsection{Q3.0-FHE-024 --- LP-013 v1\textquotesingle s "21 GPU
kernels" never had passing tests in v0.54
timeline}\label{q30-fhe-024--lp-013-v1s-21-gpu-kernels-never-had-passing-tests-in-v054-timeline}

\textbf{Severity}: Info.

\textbf{Files}: LP-013 v1 (2025-10-01), LP-013 v2 (this LP). LP-013 v2
says "v1.0 --- 21 GPU TFHE kernels." No PASS evidence on Apple M1 Max
for any of the 21. The CUDA mirror in v0.55 ("same kernel signatures,
different SM indexing") is also empty.

This is a documentation truthfulness issue, not an exploit. Flagging so
ops/marketing does not repeat the 200 tx/s/GPU-on-H100 claim on the
production landing page on Dec 25.

\begin{center}\rule{0.5\linewidth}{0.5pt}\end{center}

\subsubsection{\texorpdfstring{Q3.0-FHE-025 --- \texttt{evm256}
interpreter does not
exist}{Q3.0-FHE-025 --- evm256 interpreter does not exist}}\label{q30-fhe-025--evm256-interpreter-does-not-exist}

\textbf{Severity}: Info.

\textbf{Files}: LP-013 §"Encrypted EVM (\texttt{evm256.metal})". No file
by that name in the tree.

\textbf{Implication}: the LP-067 contract template that calls
\texttt{fhe.add} and \texttt{fhe.sub} resolves to the precompile-fhe Go
path, not to an \texttt{evm256} interpreter. The opcode-level encrypted
EVM is design-only. LP-013 v0.56 (full opcode coverage including CALL
family, LOGn) is unimplemented.

\begin{center}\rule{0.5\linewidth}{0.5pt}\end{center}

\subsection{Architectural concerns for
FHE-on-Quasar}\label{architectural-concerns-for-fhe-on-quasar}

\begin{enumerate}
\def\labelenumi{\arabic{enumi}.}
\item
  \textbf{TFHE soundness has no on-chain verification.} Quasar
  3.0\textquotesingle s cert pipeline is structured around verifying
  signatures (BLS, Ringtail, ML-DSA-Groth16) and
  \texttt{KnownTotalOrder}. An FHE attestation is a \emph{different}
  class of statement: "this circuit was evaluated correctly on this
  ciphertext under this evaluation key, producing this output
  ciphertext". This requires verifiable FHE (a SNARK over a FHE circuit)
  or a TEE attestation (LP-065). The current design pretends a HMAC over
  a hash of (circuit, input, output, BSK\_id) is enough --- it
  isn\textquotesingle t.
\item
  \textbf{F-Chain is the most centralised chain in the topology.}
  Because the secret key is held by the threshold (LP-066) or,
  currently, universally derivable (Q3.0-FHE-001), the security
  perimeter for confidentiality is the M-Chain DKG ceremony. A breach of
  the ceremony breaks every confidential balance ever encrypted. There
  is no per-tx forward secrecy; rotating the bootstrap key
  (Q3.0-FHE-008) does not re-protect old ciphertexts.
\item
  \textbf{Performance targets and reality differ by 1000× to 1.4
  million×.} The fheCRDT audit measured 0.24 LWW-Register merges/sec on
  M1 Max. LP-013 promises 200 tx/s/GPU on H100 and \textasciitilde5 tx/s
  on M1 Max. There is no kernel that delivers either number; the Go
  reference is 1,000× slower than the M1 Max claim. \textbf{This LP
  cannot be production-ready for Dec 25.}
\item
  \textbf{Quasar 3.0\textquotesingle s "one GPU process for consensus,
  DEX, EVM, and FHE in lockstep" is incompatible with Q3.0-FHE-014.} GPU
  side-channel isolation requires \emph{not} sharing the GPU between FHE
  and other workloads. Co-residency of \texttt{drain\_fhe} with
  \texttt{drain\_exec} on the same GPU (LP-013 §"Wave-tick
  co-residency") makes side-channel attacks possible from inside the EVM
  workload itself: an adversary\textquotesingle s plaintext contract can
  probe the L2 cache while a victim\textquotesingle s encrypted contract
  bootstraps. Either F-Chain runs on dedicated FHE-only GPUs, or the
  "lockstep" claim is incompatible with FHE security.
\item
  \textbf{The HMAC-keccak placeholder (LP-132 v0.38) is unsuitable for
  F-Chain at any version.} For BLS/Ringtail/ML-DSA, the placeholder is a
  temporary stand-in until v0.43--v0.45. For F-Chain, there is no "real
  verifier swap" planned, because verifiable FHE is not in any roadmap.
  \textbf{F-Chain cannot share the placeholder strategy.}
\end{enumerate}

\begin{center}\rule{0.5\linewidth}{0.5pt}\end{center}

\subsection{Must-fix-pre-launch (Dec 25, 2025) ---
ordered}\label{must-fix-pre-launch-dec-25-2025--ordered}

\begin{enumerate}
\def\labelenumi{\arabic{enumi}.}
\tightlist
\item
  \textbf{Q3.0-FHE-001} --- Replace deterministic-seed keygen with a
  real ceremony. \emph{Blocker.} Without this, FHE provides zero
  confidentiality.
\item
  \textbf{Q3.0-FHE-002} --- Remove the synchronous \texttt{decrypt}
  precompile; replace with permit-gated async threshold decryption.
\item
  \textbf{Q3.0-FHE-005} --- Implement real input ZK proof in
  \texttt{tfheVerify}; bind to caller; reject ciphertexts without proof.
\item
  \textbf{Q3.0-FHE-006} --- Implement \texttt{tfheSealOutput} with HPKE
  under the recipient\textquotesingle s public key (or remove the API).
\item
  \textbf{Q3.0-FHE-009} --- Define LP-019 §TFHE-DKG; bind
  \texttt{mchain\_ceremony\_root} to the TFHE bootstrap key transcript;
  F-Chain must verify the transcript before accepting a new
  \texttt{fchain\_fhe\_root}.
\item
  \textbf{Q3.0-FHE-003} --- Either run F-Chain inside TEEs with
  attestation via A-Chain, OR restrict to a permissioned validator set
  with explicit bond-vs-confidential-supply ratio. Pure HMAC-keccak
  F-Chain attestation must not ship.
\item
  \textbf{Q3.0-FHE-012} --- Replace underflow-revert in LP-067 transfer
  with \texttt{cmux}-based no-op-on-insufficient.
\item
  \textbf{Q3.0-FHE-015} --- Pick one \texttt{external\_product} variant;
  ban the other two on the cert path.
\item
  \textbf{Q3.0-FHE-018} --- Include destination chain ID in LP-068
  nullifier.
\item
  \textbf{Q3.0-FHE-007} --- Validate every encrypted SSTORE through the
  input proof gate before MVCC slot accepts; do not let Block-STM commit
  unverified ciphertexts.
\end{enumerate}

If any of \{001, 002, 003, 005, 009\} cannot land before Dec 25,
\textbf{F-Chain must be marked testnet-only and confidential ERC-20 /
private teleport must not be activated on mainnet.}

\subsection{Should-fix-post-launch (within 90
days)}\label{should-fix-post-launch-within-90-days}

\begin{enumerate}
\def\labelenumi{\arabic{enumi}.}
\setcounter{enumi}{10}
\tightlist
\item
  Q3.0-FHE-004 --- Ship at least the Apple Metal kernels, even if not
  optimal, so block-time targets are met.
\item
  Q3.0-FHE-008 --- Atomic key rotation with cross-key re-encryption.
\item
  Q3.0-FHE-010 --- Publish F-Chain validator set composition and bond.
\item
  Q3.0-FHE-011 --- Async decryption with nonce/permit binding.
\item
  Q3.0-FHE-013 --- Compose LP-067 with LP-064 (privacy mixer) for graph
  privacy.
\item
  Q3.0-FHE-014 --- Document GPU dedication requirement for F-Chain
  validators.
\item
  Q3.0-FHE-016 --- Property test for noise budget invariants.
\item
  Q3.0-FHE-017 --- Compiler lowering for oblivious branching.
\item
  Q3.0-FHE-019 --- Pad SSTORE gas to constant cost.
\end{enumerate}

\subsection{Lower priority}\label{lower-priority}

20--25. Q3.0-FHE-020 through 025.

\begin{center}\rule{0.5\linewidth}{0.5pt}\end{center}

\subsection{References}\label{references}

\begin{itemize}
\tightlist
\item
  LP-013 v2 (this scope): \texttt{lux/lps/LP-013-fhe-gpu.md}
\item
  LP-066: \texttt{lux/lps/LP-066-tfhe.md}
\item
  LP-067: \texttt{lux/lps/LP-067-confidential-erc20.md}
\item
  LP-068: \texttt{lux/lps/LP-068-private-teleport.md}
\item
  LP-019: \texttt{lux/lps/LP-019-threshold-mpc.md}
\item
  LP-134: \texttt{lux/lps/LP-134-lux-chain-topology.md}
\item
  LP-132 (cert lanes, HMAC placeholder):
  \texttt{lux/lps/LP-132-quasar-gpu-execution-adapter.md} lines 204-214
\item
  Go reference TFHE: \texttt{lux/fhe/}
\item
  Precompile FHE: \texttt{lux/precompile/fhe/}
\item
  Prior fhecrdt audit (corroborates noise/Solidity findings):
  \texttt{lux/audits/2026-04-12-fhecrdt-audit.tex}
\item
  Prior red-swarm audit (C-03 keygen divergence, M-03 unauth decrypt):
  \texttt{lux/audits/2026-04-13-final-red-swarm.md} lines 57-76,
  196-213.
\end{itemize}

\begin{center}\rule{0.5\linewidth}{0.5pt}\end{center}

Copyright (C) 2026, Lux Partners Limited. All rights reserved.

\end{document}
